%% sections/01-introduction.tex

\section{Introduction}
\label{sec:introduction}

Non-player characters (NPCs) in tabletop role-playing games (TRPGs) occupy
a structurally unusual position in the narrative. They are authored before
play, performed during play, and experienced by players as quasi-independent
agents with their own desires and histories. The quality of NPC interactions
is consistently cited by players as a primary driver of session
memorability~\cite{bowman2010functions, tychsen2006role}: a NPC who says
the right thing at the right moment can crystallize a session's emotional
meaning in ways that encounter design and mechanics rarely achieve.

Yet NPC design practice remains largely informal. Most NPC design advice
focuses on motivation, backstory, speech patterns, and faction
allegiance~\cite{dungeons2014players}. This advice produces plausible NPCs ---
characters who feel like people rather than combat statistics --- but it does
not reliably produce emotional payoffs. An NPC with a rich backstory may still
produce nothing memorable in play if the DM does not recognize the moment to
deploy their history, or if the party never creates the conversational opening
that would allow it.

The standard compensation for this gap is DM improvisation. Experienced DMs
develop a feel for when a party has ``earned'' an NPC beat and deliver it:
they sense when the scene is ready, judge the appropriate emotional register,
and execute the moment. This skill is real and teachable, but it is also
unreliable by design. It depends on the DM's attention state, their read of
the room, and their willingness to take a narrative risk in the middle of a
scene that may be going in a different direction. The same NPC, run by
different DMs, or by the same DM on different nights, produces inconsistent
results.

\subsection{The Arc-Completion Pattern}

This paper describes and evaluates a design pattern that addresses this
unreliability systematically: \textit{NPC arc-completion}. The pattern
has two components specified at design time, before any session begins:

\begin{itemize}
  \item \textbf{Arc-completion act:} The specific statement, act, or
    gesture the NPC is capable of making that closes their emotional arc.
    Not a range of possible responses but a specific one, written down
    verbatim or in close paraphrase before play.

  \item \textbf{Firing condition:} The exact party behavior, question,
    or action that triggers the arc-completion. Specific enough that the
    DM --- or an AI system --- knows with certainty when the condition
    has been met, without judgment calls about emotional readiness.
\end{itemize}

When the party meets the firing condition, the NPC delivers their
arc-completion by character logic: not because the DM decided it was
time, but because the pre-specified condition was satisfied. The DM
does not choose the moment; the party creates it.

This is a meaningful structural distinction. When a party earns an NPC
moment through their own actions --- when the DM's only job is to
recognize that the designed condition was met and deliver what was already
written --- the moment belongs to the session in a way that a DM-improvised
beat does not. Players may not articulate this distinction, but its effects
appear consistently in the evidence: sessions with character-logic arc-
completions score higher on narrative quality measures than sessions of
equal mechanical quality whose memorable NPC moments were DM-initiated.

\subsection{Why This Matters for AI-Simulated TRPG}

The arc-completion pattern emerged in a context that makes its benefits
particularly visible: AI-simulated TRPG, in which a large language model
(LLM) acts as Dungeon Master, runs all NPCs, and adjudicates all encounters
based on a module document and session log. In this context, there is no
``room reading'' capacity: the AI system has no access to player affect,
body language, or the social dynamics that experienced human DMs use to
time NPC moments. If the arc-completion condition is not specified in the
NPC file, the AI has no reliable basis for recognizing when the party has
earned the beat, and may either miss the moment entirely or produce a
scripted NPC speech at an inappropriate time.

The arc-completion pattern transforms this limitation into a design
requirement. Pre-writing the arc-completion forces the adventure designer
to articulate, before play, what the NPC is actually capable of saying and
under exactly what conditions. This constraint produces three advantages
over improvised NPC design: \textit{specificity} (the arc-completion text
is precise, not a paraphrase of what might be said), \textit{consistency}
(the same condition produces the same beat across multiple runs), and
\textit{DM discipline} (the system cannot extend or soften the arc-completion
because the arc-completion file specifies what happens next: nothing).

The pattern is general-purpose: while it originated in an AI-simulated
context where improvisation is architecturally unavailable, the same
template applies to human-run TRPG. In fact, the discipline requirements
of the pattern --- do not deliver early, do not extend, do not soften ---
address common failure modes in human DM improvisation that experienced
practitioners will recognize.

\subsection{Contributions}

This paper makes four primary contributions:

\begin{enumerate}
  \item \textbf{Design template:} A structured format for NPC arc-completion
    design with four fields: arc-completion act, firing condition, DM
    discipline note, and non-firing note. Illustrated with five exemplar
    NPCs drawn from the corpus (Section~\ref{sec:methodology}).

  \item \textbf{Trigger-type taxonomy:} A classification of firing condition
    types (question-answered, act-witnessed, professional-recognition,
    grief-adjacent, and situationally-triggered) with per-type firing
    rate data (Section~\ref{sec:evaluation}).

  \item \textbf{Corpus documentation:} 52 arc-completion instances across
    49 sessions and 7 complete campaigns, with session coverage rates,
    cross-archetype distribution, and simultaneous arc-convergence analysis
    (Section~\ref{sec:evaluation}).

  \item \textbf{Failure analysis:} Evidence that NPCs without arc-completion
    conditions produce memorable moments in fewer than 10\% of sessions where
    they are the primary NPC encounter, even in sessions that score 74 or
    higher on overall quality metrics (Section~\ref{sec:evaluation}).
\end{enumerate}

The pattern's primary empirical claim is that NPCs designed with explicit
arc-completion conditions fire their designed moments by character logic ---
not DM scripting --- in at least 88\% of sessions across diverse campaign
archetypes. This claim is directly testable against the session log corpus,
and Section~\ref{sec:evaluation} reports the evidence with full instance
counts, coverage breakdowns, and cross-archetype comparisons.

A secondary claim is that the pattern is thematically unconstrained: while
it was first documented in grief-adjacent campaigns, it holds equally for
morally grey NPCs (guild fences, intelligence handlers, military antagonists),
professional-code NPCs (garrison NCOs, commanding officers), and no-designed-
stakes NPCs in campaigns without a central emotional spine. Cross-archetype
validation is a core empirical contribution of this paper.

\subsection{Paper Organization}

Section~\ref{sec:related} situates the pattern in the literature on NPC design,
drama management, and character arcs in interactive narrative. Section~\ref{sec:methodology}
presents the design template, the five exemplar NPCs, and the distinction between
character-logic and DM-scripted arc moments. Section~\ref{sec:evaluation} reports
the full quantitative analysis. Section~\ref{sec:discussion} interprets the
findings and addresses limitations. Section~\ref{sec:conclusion} concludes.
